\section{零点与极值点}
\subsection{隐零点问题}
\begin{example}[2012大纲全国卷]
    已知函数 $f(x) = \frac{1}{3}x^3 + x^2 + ax$.

\begin{enumerate}
    \item 讨论 $f(x)$ 的单调性；
    \item 设 $f(x)$ 有两个极值点 $x_1, x_2$, 若过两点 $(x_1, f(x_1))$, $(x_2, f(x_2))$ 的直线 $l$ 与 $x$ 轴的交点在曲线 $y = f(x)$ 上, 求 $a$ 的值. 
\end{enumerate}
\end{example}

\begin{jie}
    \begin{enumerate}
        \item $f'(x) = x^2 + 2x + a, \Delta = 4 - 4a$.
        \begin{enumerate}
            \item 当 $\Delta = 4 - 4a \leq 0$, 即 $a \geq 1$ 时, $f'(x) \geq 0, f(x)$ 在 $\mathbb{R}$ 上是增函数；
            \item 当 $a < 1$ 时, $f'(x) = 0$, 有两个根 $x_1 = -1 - \sqrt{1-a}, x_2 = -1 + \sqrt{1-a}$. 
            \begin{itemize}
                \item 当 $x \in (-\infty, -1 - \sqrt{1-a})$ 时 $f'(x) > 0, f(x)$ 是增函数；
                \item 当 $x \in (-1 - \sqrt{1-a}, -1 + \sqrt{1-a})$ 时 $f'(x) < 0, f(x)$ 是减函数；
                \item 当 $x \in (-1 + \sqrt{1-a}, +\infty)$ 时 $f'(x) > 0, f(x)$ 是增函数. 
            \end{itemize}
        \end{enumerate}
        \item 由题设知, $x_1, x_2$ 为方程 $f'(x) = 0$ 的两个根, 故由 (1) 有 $a < 1, x_1^2 = -2x_1 - a, x_2^2 = -2x_2 - a$. 
        因此
        \begin{align*}
     f(x_1) &= \frac{1}{3}x_1^3 + x_1^2 + ax_1 = \frac{1}{3}x_1(-2x_1 - a) + x_1^2 + ax_1 \\
            &= \frac{1}{3}x_1^2 + \frac{2}{3}ax_1 = \frac{1}{3}(-2x_1 - a) + \frac{2}{3}ax_1 \\
            &= \frac{2}{3}(a-1)x_1 - \frac{a}{3}.
        \end{align*}
        同理, $f(x_2) = \dfrac{2}{3}(a-1)x_2 - \dfrac{a}{3}$. 因此直线 $l$ 的方程为 $y = \dfrac{2}{3}(a-1)x - \dfrac{a}{3}$. 
    
        设 $l$ 与 $x$ 轴的交点为 $(x_0, 0)$, 将点代入 $l$ 得, $x_0 = \dfrac{a}{2(a-1)}$. 因为点 $(x_0, 0)$ 在曲线 $y = f(x)$ 上, 故
        \begin{align*}
          f(x_0) &= \frac{1}{3}\left[\frac{a}{2(a-1)}\right]^3 + \left[\frac{a}{2(a-1)}\right]^2 + \frac{a^2}{2(a-1)} \\
            &= \frac{a^2}{24(a-1)^3}(12a^2 - 17a + 6) = 0,
        \end{align*}
        解得 $a = 0$ 或 $a = \dfrac{2}{3}$ 或 $a = \dfrac{3}{4}$. 
    \end{enumerate}
\qed
\end{jie}

\begin{example}[2012课标全国卷]
    设函数 $f(x) = e^x - ax - 2$.

\begin{enumerate}
    \item 求 $f(x)$ 的单调区间；
    \item 若 $a = 1, k$ 为整数, 且当 $x > 0$ 时, $(x-k)f'(x) + x + 1 > 0$, 求 $k$ 的最大值. 
\end{enumerate}
\end{example}

\begin{jie}
\ 


\((1)\)$f(x)$ 的定义域为 $(-\infty, +\infty)$ $f'(x) = e^x - a$. 
        \begin{itemize}
            \item 若 $a \leq 0$, 则 $f'(x) > 0$, 所以 $f(x)$ 在 $(-\infty, +\infty)$ 上单调递增. 
            \item 若 $a > 0$, 令 $f'(x) = e^x - a = 0$, 得 $x = \ln a$. 
            \begin{itemize}
                \item 当 $x \in (-\infty, \ln a)$ 时 $f'(x) < 0$；当 $x \in (\ln a, +\infty)$ 时 $f'(x) > 0$. 所以 $f(x)$ 在 $(-\infty, \ln a)$ 上单调递减, 在 $(\ln a, +\infty)$ 上单调递增. 
            \end{itemize}
        \end{itemize}
        综上所述, 
        \begin{itemize}
            \item 当 $a \leq 0$ 时 $f(x)$ 的增区间为 $(-\infty, +\infty)$；
            \item 当 $a > 0$ 时 $f(x)$ 的减区间为 $(-\infty, \ln a)$, 增区间为 $(\ln a, +\infty)$. 
        \end{itemize}

\((2)\)  
由于 $a = 1$, 由 (1) 可得, $(x-k)f'(x) + x + 1 = (x-k)(e^x - 1) + x + 1$. 
故当 $x > 0$ 时, $(x-k)f'(x) + x + 1 > 0$ 等价于 $$k < \frac{x+1}{e^x - 1} + x (x > 0)$$. 

 令 $g(x) = \dfrac{x+1}{e^x - 1} + x$, 则 $g'(x) = \dfrac{-xe^x - 1}{(e^x - 1)^2} + 1 = \dfrac{e^x(e^x - x - 2)}{(e^x - 1)^2}$. 
 由于\(e^x>0\)恒成立, 仅需要研究\(h(x) = e^x - x - 2\)的零点即可. 
 
 由 (1) 知, 函数 $h(x) = e^x - x - 2$ 在 $(0, +\infty)$ 上单调递增. 
而 $h(1) = e - 3 < 0, h(2) = e^2 - 4 > 0$, 所以 $h(x)$ 在 $(0, +\infty)$ 上存在唯一的零点, 故 $g'(x)$ 在 $(0, +\infty)$ 上存在唯一的零点. 
 设此零点为 $x_0$, 则 $x_0 \in (1, 2)$. 

当 $x \in (0, x_0)$ 时, $g'(x) < 0, g(x)$ 为减函数；当 $x \in (x_0, +\infty)$ 时, $g'(x) > 0, g(x)$ 为增函数. 
 所以 $g(x)$ 在 $(0, +\infty)$ 上的最小值为 $g(x_0)$. 
 又由 $g'(x_0) = 0$, 可得 $e^{x_0} = x_0 + 2$, 所以 $g(x_0) = x_0 + 1 \in (2, 3)$. 
由于等价于 $k < g(x_0)$, 故整数 $k$ 的最大值为 2. 
\qed
\end{jie}
\begin{example}[2017年II卷理]
    已知函数 $f(x) = x^2 - x - x\ln x$, 证明：$f(x)$ 存在唯一的极大值点 $x_0$, 且 $e^{-2} < f(x_0) < 2^{-2}$. 
\end{example}
\begin{jie}
    首先, 计算得
    \[
    f'(x) = 2x - 2 - \ln x, \quad f''(x) = \frac{2x - 1}{x},
    \]
    因此 $f'(x)$ 在 $\left(0, \frac{1}{2}\right)$ 内单调递减, 在 $\left(\frac{1}{2}, +\infty\right)$ 内单调递增. 又
    \[
    f'\left(\frac{1}{e^2}\right) = \frac{2}{e^2} > 0, \quad f'\left(\frac{1}{2}\right) = \ln 2 - 1 < 0, \quad f'(1) = 0,
    \]
    因此存在 $x_0 \in \left(\frac{1}{e^2}, \frac{1}{2}\right)$, 使得 $f'(x_0) = 0$. 
    
    在上面我们得到了 $f'(x_0) = 2x_0 - 2 - \ln x_0 = 0$. $f(x)$ 在 $(0, x_0)$ 内单调递增, 在 $(x_0, 1)$ 内单调递减, 在 $(1, +\infty)$ 内单调递增, 因此 $x = x_0$ 是 $f(x)$ 的极大值点. 一方面, 有
    \[
    f(x_0) = x_0^2 - x_0 - x_0\ln x_0 = x_0^2 - x_0 - x_0(2x_0 - 2) = x_0(1 - x_0) < \frac{1}{4} = 2^{-2};
    \]
    另一方面, 根据 $x_0$ 是极大值点得
    \[
    f(x_0) > f\left(\frac{1}{e}\right) = e^{-2},
    \]
    因此 $f(x)$ 存在唯一的极大值点 $x_0$, 且 $e^{-2} < f(x_0) < 2^{-2}$. 








    \qed
\end{jie}
\begin{exercise}[2019年II卷文数]
    已知函数 $f(x) = (x-1)\ln x - x - 1$. 证明：
\begin{enumerate}
    \item $f(x)$ 存在唯一的极值点；
    \item $f(x) = 0$ 有且仅有两个实根, 且两个实根互为倒数. 
\end{enumerate}

\end{exercise}

\begin{proof}
\ 

\((1)\)$f(x)$ 的定义域为 $(0, +\infty)$, 计算得 $f'(x) = \ln x - \dfrac{1}{x}$, 由 $\ln x$ 与 $-\dfrac{1}{x}$ 在 $(0, +\infty)$ 内单调递增知 $f'(x)$ 在 $(0, +\infty)$内单调递增, 计算得
        \[ f'(1) = -1 < 0,  \quad f'(e) = 1 - \frac{1}{e} > 0. \]
     故根据零点定理知, 存在 $x_0 \in (1, e)$, 使得 $f'(x_0) = 0$, $f(x)$ 在 $(0, x_0)$ 内单调递减, 在 $(x_0, +\infty)$ 内单调递增, $x = x_0$ 是 $f(x)$
     的极小值点. 

\((2)\)计算得 $f(e) = -2 < 0$, $f(e^2) = e^2 - 3 > 0$, 故根据零点定理知, 存在 $x_1 \in (e, e^2)$, 使得 $f(x_1) = 0$
    . 注意到此时
        \[
        f
    \left(\frac{1}{x}\right) = -\left(\frac{1}{x} - 1\right)\ln x - \frac{1}{x} - 1 = \frac{(x-1)\ln x - x - 1}{x} = \frac
    {f(x)}{x},
        \]
        因此当且仅当 
    $f\left(\frac{1}{x}\right) = 0$ 时 $f(x) = 0$, 从而 $x_2 = \frac{1}{x_1}$ 也是 $f(x)$ 的零点, 根据单调性知 $f(x)$ 有且仅有两个零点 $x_1$ 和 $x_2$
    . 

    \qed
\end{proof}

\begin{example}[2022年新高考I卷]
    已知函数\(f(x)=e^x-ax\)和\(g(x)=ax-\ln x\)有相同的最小值.
    \begin{enumerate}
        \item 求\(a\)的值；
        \item 证明：存在直线 $y = b$, 其与两条曲线 $y = f(x)$ 和 $y = g(x)$ 共有三个不同的交点, 并且从左到右的三个交点的横坐标成等差数列. 
    \end{enumerate}
\end{example}
\begin{jie}
\  

\((1)\)$f(x)$ 的定义域为 $\mathbb{R}$, $f'(x) = e^x - a$, 由 $f(x)$ 有最小值得 $a > 0$, 且 $f(x)_{\min} = f(\ln a) = a - a\ln a$；

同时\(g(x)\)的定义域为\((0,+\infty)\), $g'(x) = \dfrac{ax - 1}{x}$, 
从而 $g(x)_{\min} = g\left(\dfrac{1}{a}\right) = 1 + \ln a$. 

由\(f_{\min}(x)=g_{\min}(x)\)得
$$ \ln a - \dfrac{a-1}{a+1} = 0$$. 
            
        令 $h(a) = \ln a - \dfrac{a-1}{a+1}$, 其中 $a > 0$, 计算得
        \[ h'(a) = \frac{1}{a} - \frac{2}{(a+1)^2} = \frac{a^2 + 1}{a(a+1)^2} > 0, \]
        因此 $h(a)$ 在 $(0, +\infty)$ 内单调递增, 注意到 $h(1) = 0$, 于是 $h(a)$ 的唯一零点是 $a = 1$. 
        
\((2)\)不妨设 $x_1 < 0 < x_2 < 1 < x_3$ , 此时
    \[e^{x_1} - x_1 = e^{x_2} - x_2 = x_2 - \ln x_2 = x_3 - \ln x_3. \]

    首先, 注意到
        \[f(x_1) = e^{x_1} - x_1 = g(x_2) = x_2 - \ln x_2 = e^{\ln x_2} - \ln x_2 = f(\ln x_2),\]
    其中 $x_1 < 0, \ln x_2 < 0$, 根据 $f(x)$ 在 $(-\infty, 0)$ 内单调递减, 
    知 $x_1 = \ln x_2$, 即 $x_2 = e^{x_1}$. 
    
    接下来, 注意到
    \[f(x_2) = e^{x_2} - x_2 = g(x_3) = x_3 - \ln x_3 = e^{\ln x_3} - \ln x_3 = f(\ln x_3),  \]
    其中 $x_2 > 0, \ln x_3 > 0$, 根据 $f(x)$ 在 $(0, +\infty)$ 内单调递增, 知 $x_2 = \ln x_3$, 即 $x_3 = e^{x_2}$. 最后, 计算得
     \[x_1 + x_3 = \ln x_2 + e^{x_2} = 2x_2,\]
    从而 $x_1, x_2, x_3$ 成等差数列. 
    \qed
\end{jie}



含参的处理：

\begin{example}[2009 全国 II 理 ]
    设 $f(x) = x^2 + a \ln(1 + x)$ 有两个极值点 $x_1, x_2$, 且 $x_1 < x_2$. 
\begin{enumerate}
    \item 求实数 $a$ 的取值范围, 并讨论 $f(x)$ 的单调性；
    \item 证明：$f(x_2) > \dfrac{1 - 2\ln 2}{4}$. 
\end{enumerate}
\end{example}

\begin{jie}
    \                  

    \((1)\)原问题等价于 $f'(x) = \dfrac{2x^2 + 2x + a}{1 + x}$ 在 $(-1, +\infty)$ 上有两个零点 $x_1, x_2$, 即 $h(x) = 2x^2 + 2x + a$ 在 $(-1, +\infty)$ 上有两个零点 $x_1, x_2$, 如图 5-22 所示, 则原问题等价于
    \[
    \begin{cases}
        \Delta = 4 - 8a > 0 \\
        h(-1) > 0
    \end{cases}
    \]
    解得 $0 < a < \dfrac{1}{2}$, 且此时 $f(x)$ 在 $(-1, x_1)$ 上单调递增, 在 $(x_1, x_2)$ 上单调递减, 在 $(x_2, +\infty)$ 上单调递增. 

    \((2)\)原问题等价于证明
    \[
    \begin{cases}
        2x_2^2 + 2x_2 + a = 0 \\
        f(x_2) = x_2^2 + a \ln(1 + x_2) > \dfrac{1 - 2\ln 2}{4}
    \end{cases}
    \]
    在 $a \in \left(0, \frac{1}{2}\right)$上恒成立. 由第一个方程可得
    \[a = -2x_2^2 - 2x_2 \]
    代入式 (5.6.1) 的第二个方程可得
    \[f(x_2) = x_2^2 - 2(x_2 + x_2^2) \ln(1 + x_2) > \frac{1 - 2\ln 2}{4} \]
    显然不等式并不是对所有的 $x_2 \in (-1, +\infty)$ 都成立（$x_2$ 趋于 $+\infty$ 会产生矛盾）. 
    考虑\(h(0)=a\in(0,\dfrac{1}{2})\), 可知 $x_2 \in \left(-\dfrac{1}{2}, 0\right)$. 

    故原问题等价于证明 $$g(t) = t^2 - 2(t + t^2) \ln(1 + t) > \frac{1 - 2\ln 2}{4}$$
    在 $\left(-\dfrac{1}{2}, 0\right)$上恒成立, 求导得
    \[g'(t) = -2(1 + 2t) \ln(1 + t) > 0, \]
    即 $g(t)$ 在 $\left(-\dfrac{1}{2}, 0\right)$ 上单调递增. 
    故 $g(t) > g\left(-\dfrac{1}{2}\right) = \dfrac{1}{4} - \dfrac{2\ln 2}{4}$, 这就证明了 $f(x_2) > \dfrac{1 - 2\ln 2}{4}$. 

    \qed
\end{jie}

\begin{example}[2021年天津]
    已知 $a > 0$, 函数 $f(x) = ax - xe^x$. 
\begin{enumerate}
    \item 证明 $f(x)$ 存在唯一的极值点；
    \item 若存在 $a$, 使得 $f(x) \leqslant a + b$ 对任意 $x \in \mathbb{R}$ 成立, 求实数 $b$ 的取值范围. 
\end{enumerate}
\end{example}
\begin{jie}
    \ 

    \((1)\)计算得
    \[
    f'(x) = a - (x + 1)e^x, \quad f''(x) = -(x + 2)e^x,
    \]
    因此 $f'(x)$ 在 $(-\infty, -2)$ 内单调递增, 在 $(-2, +\infty)$ 内单调递减, 
    当 $x \to -\infty$ 时, $f'(x) \to a > 0$, 且 $f'(-1) = a > 0$, 
    当 $x \to +\infty$ 时, $f'(x) \to -\infty$, 
    因此存在 $x_0 \in (-1, +\infty)$, 使得 $f'(x_0) = 0$, 
    $x = x_0$ 是 $f(x)$ 的极值点, 且满足 $(x_0 + 1)e^{x_0} = a$. 


    \((2)\)
    首先固定 $a$, 则由 $(x_0 + 1)e^{x_0} = a$ 知 $x_0$ 也固定. 令 $g(x) = f(x) - a$, 则 $g'(x) = f'(x)$, $g(x)$ 在 $(-\infty, x_0)$ 内单调递增, 在 $(x_0, +\infty)$ 内单调递减, 从而
    \[g(x)_{\max} = g(x_0) = ax_0 - x_0e^{x_0} - a = (x_0^2 - x_0 - 1)e^{x_0} \leqslant b.\]
    令 $\varphi(x) = (x^2 - x - 1)e^x$, 其中 $x > -1$, 计算得
    \[\varphi'(x) = (x^2 + x - 2)e^x = (x - 1)(x + 2)e^x,\]
    因此 $\varphi(x)$ 在 $(-1, 1)$ 内单调递减, 在 $(1, +\infty)$ 内单调递增, 
    $\varphi(x)_{\max} = \varphi(1) = -e$, 因此实数 $b$ 的取值范围是 $[-e, +\infty)$. 

    \qed
\end{jie}

















\begin{example}
    求证：\(e^x-\ln x >2\).
\end{example}




















\begin{example}
    求证：\(e^x-\ln x >2.3\);（参考值：\(\sqrt{\e}=1.649,\ln 2=0.693,\dfrac{4}{3}\sqrt{3}=2.3094\)）.
\end{example}

\begin{proof}
    设 $f(x) = e^x - \ln x$, 定义域 $x \in (0, +\infty)$, 则 $f'(x) = e^x - \frac{1}{x}$ 在 $x \in (0, +\infty)$ 上为增函数. 

又 $f'\left(\frac{1}{2}\right) = \sqrt{e} - 2 < 0$, $f'(1) = e - 1 > 0$, 所以存在唯一 $x_0 \in \left(\frac{1}{2}, 1\right)$, 使得 $f'(x_0) = 0$. 

所以当 $x \in (0, x_0)$ 时 $f'(x) < 0$, $f(x)$ 在 $x \in (0, x_0)$ 上单调递减；

当 $x \in (x_0, +\infty)$ 时 $f'(x) > 0$, $f(x)$ 在 $x \in (x_0, +\infty)$ 上单调递增. 

所以 $f(x)_{\min} = f(x_0) = e^{x_0} - \ln x_0$, 

又 $f'(x_0) = 0$, 即 $e^{x_0} = \frac{1}{x_0}$, 两边取对数得 $x_0 = -\ln x_0$, 都代入得

\[ f(x)_{\min} = f(x_0) = \frac{1}{x_0} + x_0 \], 在 $x_0 \in \left(\frac{1}{2}, 1\right)$ 上为减函数, 

所以 $f(x_0) > f(1) = 2$, 不满足题意. 

又 $f\left(\frac{1}{2}\right) = \sqrt{e} + \ln 2 > 2.3$, 说明 $x_0$ 范围估值不够精确, 需进一步缩小范围. 

又 $x_0 = -\ln x_0$ 说明 $x_0$ 是一个不精确计算出的值, 

而当 $x \in (0, 1)$ 时, 有 $\ln x > \frac{1}{2}(x - \frac{1}{x})$, 于是

\[ x = -\ln x < -\frac{1}{2}(x - \frac{1}{x}) \], 解得 $x < \frac{1}{\sqrt{3}}$, 

所以进一步确定 $x_0 \in \left(\dfrac{1}{2}, \dfrac{1}{\sqrt{3}}\right)$, 有 $f(x_0) > f\left(\dfrac{1}{\sqrt{3}}\right) = \sqrt{3} + \dfrac{1}{\sqrt{3}} = \dfrac{4\sqrt{3}}{3} > 2.3$. 

综上所述, $e^x - \ln x > 2.3$. 
\qed
\end{proof}

\begin{definition}
    设函数\(f(x)=x\e ^x,x\in\mathbb{R}\), 则记\(W(x)\)为\(f(x)(x\geq -1)\)的反函数, \(W_{-1}(x)\)为\(f(x)(x\leq -1)\)
    的反函数. 函数\(W(x)\)和\(W_{-1}(x)\)称为Lambert函数. (Lambert W Functionf)
\end{definition}




\begin{example}[25.7第五届久洵杯19]
    已知函数 $f(x) = \ln \frac{1+x}{1-x} - a$, \(g(x) = 2x^2 - a\sin x$. 
\begin{enumerate}
    \item 讨论 $f(x)$ 的单调性；
    \item 若 $g(x)$ 在 $\left(-\frac{\pi}{2}, \frac{\pi}{2}\right)$ 有且只有 2 个零点, 求 $a$ 的取值范围；
    \item 设 $a > 0$. 在 2. 的条件下, 记 $f(x)$ 的零点为 $x_1$,
     $g(x)$ 在 $\left(-\frac{\pi}{2}, \frac{\pi}{2}\right)$ 的零点为 $x_2$ $(x_2 \neq 0)$, 
    证明：$x_1 < \sin x_2$. 
\end{enumerate}

\end{example}

\begin{jie}
    
\end{jie}

\clearpage
\subsection{多变量的基本处理}
\clearpage
\subsection{极值点偏移}

\begin{example}[2010年天津理]
    已知函数 $f(x) = \dfrac{x}{e^x}$, 如果 $x_1 \neq x_2$, 且 $f(x_1) = f(x_2)$, 证明：$x_1 + x_2 > 2$. 
\end{example}

\begin{proof}
    $f'(x) = (1-x)e^{-x}$, 故 $f(x)$ 在 $(-\infty, 1)$ 上递增, 在 $(1, +\infty)$ 上递减. 不妨设 $x_1 < x_2$, 由 $f(x_1) = f(x_2)$ 可得 $x_1 < 1$, $x_2 > 1$, 由 $f(x_1) = f(x_2)$ 可得
    \[
    x_1 + x_2 > 2 \Rightarrow x_1 > 2 - x_2 \Rightarrow f(x_1) > f(2 - x_2) \Rightarrow f(x_2) > f(2 - x_2)
    \]
    即只需证 $g(x) = f(x) - f(2-x) = xe^{-x} - (2-x)e^{x-2} > 0$ 在 $(1, +\infty)$ 上成立. 求导可得 $g'(x) = (x-1)(e^{x-2} - e^{-x})$. 因为 $h(x) = e^{x-2} - e^{-x}$ 在 $(1, +\infty)$ 上递增, 所以 $h(x) > h(1) = 0$, 因此 $g'(x) = (x-1)h(x) > 0$ 在 $(1, +\infty)$ 上恒成立. 因此 $g(x)$ 在 $(1, +\infty)$ 上递增, 故 $g(x) > g(1) = f(1) - f(1) = 0$, 即 $f(x) > f(2-x)$, 因此 $f(x_2) > f(2-x_2)$. 故
    \[
    f(x_2) > f(2 - x_2) \Rightarrow f(x_1) > f(2 - x_2) \Rightarrow x_1 + x_2 > 2
    \]
    
\end{proof}


\begin{example}
    若 $m > 1$ 且 $x_1, x_2$ 为 $g(x) = \ln(x+m) - x$ 的两个零点, 证明：$x_1 + x_2 < 0$. 
\end{example}
\begin{analyse}
    \(0\)不是原函数的极值点. 考虑分参. 
\end{analyse}
\begin{proof}
    因为 $x_1, x_2$ 为 $g(x)$ 的两个零点, 所以 $x_1, x_2$ 为 $h(x) = e^x - x - m$ 的两个零点. 求导可得 $h'(x) = e^x - 1$, 故 $h(x)$ 在 $(0, +\infty)$ 上单调递增, 在 $(-m, 0)$ 上单调递减. 不妨设 $x_1 < x_2$, 由 $h(x_1) = h(x_2)$ 可知 $x_1 < 0 < x_2$
, 则
\[
x_1 + x_2 < 0 
\Leftrightarrow x_1 < -x_2 \Leftrightarrow h(x_1) > h(-x_2) \Leftrightarrow
 h(x_2) > h(-x_2)
\]
然后构造函数 
$F(x) = h(x) - h(-x)$, 因为 $h(x)$ 的定义域为 $(-m, +\infty)$, 所以 $h(-x)$ 的定义域为 $(-\infty, m)$, 因此 $F(x)$ 的定义域为 $(-m, m)$. 又因为我们关心的是 $F(x_2)$ 的符号且 $x_2 > 0$, 所以只需证明函数 $F(x) > 0$ 在 $(0, m)$ 上成立即可. 因为 $F(0) = 0$, 所以只需证明 $F(x)$ 在 $(0, m)$
 上单调递增即可. 

因为 
$F(x) = e^x - \frac{1}{e^x} - 2x$, 则求导可得 $F'(x) = e^x + \frac{1}{e^x} - 2$, 当 $x \in (0, m)$ 时, $F'(x) > 0$, 即 $F(x)$ 在 $(0, m)$ 上单调递增, 故 $x_1 + x_2 < 0$
 得证. 

\end{proof}

\begin{example}
    已知函数 $f(x) = \frac{1 + \ln x}{x}$, 若 $f(x) = m$ 有两个不同实根 $x_1, x_2$, 求证：$\frac{1}{x_1} + \frac{1}{x_2} > 2$. 
\end{example}

\begin{proof}
    所证结论 $\frac{1}{x_1} + \frac{1}{x_2} > 2$ 并不是常见类型, 怎么办？把它化成常规类型, 通常情况下都是换元法：令 $t_1 = \frac{1}{x_1}, t_2 = \frac{1}{x_2}$, 即证明 $t_1 + t_2 > 2$. 接下来需要找出 $t_1, t_2$ 所满足的方程. 

由 $f(x_1) = f(x_2) = m$, 可知 $1 + \ln x_1 = mx_1$, $1 + \ln x_2 = mx_2$, 可得 $t_1 - t_1 \ln t_1 = t_2 - t_2 \ln t_2$. 令 $g(t) = t - t \ln t$, 则 $g(t_1) = g(t_2)$. 

则原问题等价于“已知函数 $g(t) = t - t \ln t$, 且 $t_1, t_2$ 满足 $g(t_1) = g(t_2)$, 证明：$t_1 + t_2 > 2$”. 求导可得 $g'(t) = -\ln t$, 则 $g(t)$ 在 $(1, +\infty)$ 上单调递减, 在 $(0, 1)$ 上单调递增, 因为 $g(t_1) = g(t_2)$, 不妨设 $0 < t_1 < 1 < t_2$, 则
\[
t_1 + t_2 > 2 \Leftrightarrow t_1 > 2 - t_2 \Leftrightarrow g(t_1) > g(2 - t_2) \Leftrightarrow g(t_2) > g(2 - t_2)
\]
求导可得然后构造辅助函数 $F(t) = g(t) - g(2-t)$, 因为 $g(t)$ 的定义域为 $(0, +\infty)$, 所以 $g(2-t)$ 的定义域为 $(-\infty, 2)$, 因此 $F(t)$ 的定义域为 $(0, 2)$. 又因为我们关心的是 $F(t_2)$ 的符号且 $t_2 > 1$, 所以只需证明函数 $F(t) > 0$ 在 $(1, 2)$ 上成立即可. 因为 $F(1) = 0$, 所以只需证明 $F(t)$ 在 $(1, 2)$ 上单调递增即可. 

令 $F(t) = g(t) - g(2-t)$, $t \in (1, 2)$, 则 $F'(t) = -\ln(2-t) - \ln t = -\ln t(2-t) > -\ln\left(\frac{t+2-t}{2}\right)^2 = 0$, 即 $F(t)$ 在 $(1, 2)$ 上单调递增. 故 $\frac{1}{x_1} + \frac{1}{x_2} > 2$ 得证. 
\end{proof}
\begin{example}
    若 $f(x) = x^2 \ln x - m$ 有两个不同零点 $x_1, x_2$, 求证：$x_1^2 + x_2^2 > \frac{2}{e}$. 
\end{example}
\begin{proof}
    令 $t_1 = x_1^2, t_2 = x_2^2$, 由 $f(x_1) = f(x_2)$, 可得 $x_1^2 \ln x_1 = x_2^2 \ln x_2$, 即 $t_1 \ln t_1 = t_2 \ln t_2$, 记函数 $h(t) = t \ln t$, 即证 $t_1 + t_2 > \frac{2}{e}$. 求导可得 $h'(t) = 1 + \ln t$, 当 $x \in \left(0, e^{-1}\right)$ 时, 函数 $h(x)$ 单调递减, 当 $x \in \left(e^{-1}, +\infty\right)$ 单调递增. 不妨设 $t_1 < t_2$, 因为 $h(t_1) = h(t_2)$, 所以 $t_1 < \frac{1}{e} < t_2$, 则
\[
t_1 + t_2 > \frac{2}{e} \Leftrightarrow t_1 > \frac{2}{e} - t_2 \Leftrightarrow h(t_1) > h\left(\frac{2}{e} - t_2\right) \Leftrightarrow h(t_2) > h\left(\frac{2}{e} - t_2\right)
\]
令 $F(t) = h(t) - h\left(2e^{-1} - t\right)$, 因为 $h(t)$ 的定义域为 $(0, +\infty)$, 所以 $h\left(2e^{-1} - t\right)$ 的定义域为 $(-\infty, 2e^{-1})$. $F(t) = h(t) - h\left(2e^{-1} - t\right)$, 因为 $h(t)$ 的定义域为 $(0, 2e^{-1})$ 上的符号, 因此只考虑 $F(t)$ 在 $(e^{-1}, 2e^{-1})$ 上的符号即可. 求导可得 $F'(t) = 2 + \ln\left(t(2e^{-1} - t)\right) < 0$, 故 $F(t)$ 在 $(e^{-1}, 2e^{-1})$ 上递减, 即 $F(t_2) < F(e^{-1}) = 0$, 故 $x_1^2 + x_2^2 > \frac{2}{e}$ 得证. 
\end{proof}
\begin{example}
    已知函数 $f(x) = \ln x - ax$, 若函数 $f(x)$ 有两个零点 $x_1, x_2$, 证明：$x_1 x_2 > e^2$. 
\end{example}

\begin{proof}
    显然 $x = e$ 不为 $f(x)$ 的极值点, 因此先通过分离变量可得 $a = \frac{\ln x}{x}$. 令 $g(x) = \frac{\ln x}{x}$. 求导可得 $g'(x) = \frac{1 - \ln x}{x^2}$, 则 $g(x)$ 在 $(0, e)$ 上递增, 在 $(e, +\infty)$ 上递减. 不妨设 $0 < x_1 < x_2$ 且 $g(x_1) = g(x_2)$, 则 $0 < x_1 < e < x_2$, 故
\[
x_1 x_2 > e^2 \Leftrightarrow x_1 > \frac{e^2}{x_2} \Leftrightarrow g(x_1) > g\left(\frac{e^2}{x_2}\right) \Leftrightarrow g(x_2) > g\left(\frac{e^2}{x_2}\right)
\]
令 $h(x) = g(x) - g\left(\frac{e^2}{x}\right)$, 因为 $g(x)$ 的定义域为 $(0, +\infty)$, 所以 $g\left(\frac{e^2}{x}\right)$ 的定义域为 $(0, +\infty)$, 故 $h(x)$ 的定义域为 $(0, +\infty)$. 又因为我们关心的是 $h(x_2)$ 的符号, 因此我们只需研究 $h(x)$ 在 $(e, +\infty)$ 上的符号即可. 求导可得 $h'(x) = (1 - \ln x)\left(\frac{1}{x^2} - \frac{1}{e^2}\right) > 0$ 在 $(e, +\infty)$ 上恒成立. 故 $h(x_2) > h(e) = 0$, 即 $x_1 x_2 > e^2$ 得证. 
\end{proof}

\begin{example}[2021年新高考I卷]
    已知函数 $f(x) = x(1 - \ln x)$. 

\begin{enumerate}
    \item 讨论 $f(x)$ 的单调性；
    \item 设 $a, b$ 为两个不相等的正实数, 且 $b \ln a - a \ln b = a - b$, 证明：$2 < \frac{1}{a} + \frac{1}{b} < e$. 
\end{enumerate}

\end{example}
\begin{jie}
    
\end{jie}


极值点被放缩过

\begin{example}
    已知函数 $f(x) = \frac{1}{2}x^2 + (1-a)x - a\ln x, a \in \mathbb{R}$.
若 $f(x)$ 存在两个不同零点 $x_1, x_2$, 求证：$x_1 + x_2 > 2$. 
\end{example}
\begin{jie}
    \begin{align*}
        f'(x) &= x + 1 - a - \frac{a}{x} = (x + 1) - \frac{a(x + 1)}{x} \\
        &= (x + 1)\left(1 - \frac{a}{x}\right) \\
        \text{if } a \leqslant 0, & \ f'(x) > 0, \text{ wrong} \Rightarrow a > 0 \\
        x \in (0, a), & \ f'(x) < 0, x \in (a, +\infty), f'(x) > 0 \\
        f'(x) = 0 & \Rightarrow x = a
        \end{align*}
        
        我们发现, 这里的极值点是 $a$, 如果是标准的偏移, 应该是证明 $x_1 + x_2 > 2a$. 但题目却嚷证明 $x_1 + x_2 > 2$, 我们一时难以理解. 
        \begin{align*}
            f(x)_{\min} &= f(a) = \frac{a^2}{2} + (1-a)a - a\ln a < 0 \\
            \Rightarrow \frac{a}{2} + 1 - a - \ln a &= 1 - \frac{a}{2} - \ln a < 0 \\
            h(a) &= 1 - \frac{a}{2} - \ln a, h(1) = 1 - \frac{1}{2} = \frac{1}{2} > 0 \\
            \Rightarrow a &> 1
            \end{align*}
            
            要使 $f(x)$ 有两个不同零点, 则 $f(x)$ 的最小值（极小值）必小于零, 因此, 我们要解出这个关于 $a$ 的不等式的解集. 然而, 我们发现这个不等式是超越的. 
            
            虽然不能解出 $a$ 的具体范围, 但我们可以得到一个必要条件, 即 $a > 1$, 所以, 要证明原问题的 $x_1 + x_2 > 2$, 我们只需要证明 $x_1 + x_2 > 2a$ 即可. 而 $2a$ 刚好是我们 2 倍的极值点, 所以, 我们已经把问题转化成了极值点偏移的标准形式. 
            
            \begin{align*}
            f'(x) &= x + 1 - a - \frac{a}{x} = (x + 1) - \frac{a(x + 1)}{x} \\
            &= (x + 1)\left(1 - \frac{a}{x}\right) \\
            \text{if } a \leqslant 0, & \ f'(x) > 0, \text{ wrong} \Rightarrow a > 0 \\
            x \in (0, a), & \ f'(x) < 0, x \in (a, +\infty), f'(x) > 0 \\
            f'(x) = 0 & \Rightarrow x = a \\
            f(x)_{\min} &= f(a) = \frac{a^2}{2} + (1-a)a - a\ln a < 0 \\
            \Rightarrow \frac{a}{2} + 1 - a - \ln a &= 1 - \frac{a}{2} - \ln a < 0 \\
            h(a) &= 1 - \frac{a}{2} - \ln a, h(1) = 1 - \frac{1}{2} = \frac{1}{2} > 0 \\
            \Rightarrow a &> 1, 0 < x_1 < a < x_2, 2a - x_2 < a \\
            x_1 + x_2 > 2a &\Leftarrow x_1 > 2a - x_2 \Leftarrow f(x_1) < f(2a - x_2) \\
            \Leftrightarrow f(x_2) - f(2a - x_2) &< 0, \\
            F(x) &= f(x) - f(2a - x) = 2x - a\ln x + a\ln(2a - x) - 2a \\
            F'(x) &= 2 - \frac{a}{x} - \frac{a}{2a - x} = -\frac{2(x - a)^2}{x(2a - x)} \leqslant 0 \\
            \therefore F(x) &< F(a) = 0 \Rightarrow x_1 + x_2 > 2a > 2
            \end{align*}
            
            \qed
\end{jie}





\clearpage
\subsection{零点差}
\clearpage
